\section{Introduction} \label{s0}

It is well-known that starting with the famous Lichnerowicz vanishing theorem~\cite{L63}, Dirac operators have played important roles in the study of Riemannian metrics of positive scalar curvature on spin manifolds (cf.~\cite{GL83} and~\cite{LaMi89}). A notable example is Llarull's rigidity theorem~\cite{L98} which states that for a compact spin Riemannian manifold $\big(M,g^{TM}\big)$ of dimension $n$ such that the associated scalar curvature $k^{TM}$ verifies that $k^{TM}\geq n(n-1)$, any (non-strict) area decreasing smooth map $f\colon M\rightarrow S^n(1)$ of nonzero degree is an isometry.

Recently, Gromov states in \cite[p.~45]{Gr19} a noncompact extension of the Llarull theorem. Namely, if $\big(M,g^{TM}\big)$ is an $n$ dimensional noncompact complete spin Riemannian manifold, $f\colon M\rightarrow S^n(1)$ a smooth (non-strict) area decreasing map (which is locally constant near infinity) of nonzero degree such that $k^{TM}\geq n(n-1)$ on the support of~${\rm d}f$, then
\begin{gather}\label{0.1}
\inf \big(k^{TM}\big)\leq 0.
\end{gather}

The argument used by Gromov for~(\ref{0.1}) relies on the relative index theorem of Gromov--Lawson~\cite{GL83}, which depends on the positivity of $k^{TM}$ near infinity. Gromov then raises the question that whether the inequality in (\ref{0.1}) can actually be made {strict}.

The purpose of this short note is to provide a positive answer to this question when~$n$ is even. That is, (\ref{0.1}) can indeed be improved to $\inf \big(k^{TM}\big)< 0$. When $n$ is odd, we improve~(\ref{0.1}) to $\inf \big(k^{TM}\big)< 0$ under the condition that $k^{TM}> n(n-1)$ on the support of~${\rm d}f$.

The main idea of the proof, similar to \cite[equation~(1.11)]{Z19}, is to deform the involved twisted Dirac operator (constructed as in~\cite{L98})
on~$M$ by a suitable endomorphism of the twisted vector bundle (cf.~(\ref{1.7})). The deformed Dirac operator turns out to be invertible near infinity, and one can then apply the relative index theorem to complete the proof.

\section{The main results and their proofs}\label{s1}
